\input _template

\setlanguage{CS}

\Title{Průměry}

\Author{Káťa Danilina}

\sec Úvod

Představme si častou situaci: z~předmětu ve škole máme následující známky
$$
1,\ 1,\ 3,\ 2,\ 1
$$
a chceme zjistit, co nejspíš dostaneme na vysvědčení. Když si spočteme
$$
\frac{1+1+3+2+1}{5}=1{,}6,
$$
tak se to naneštěstí pro nás zaokrouhlí na~2. My si ale ukážeme, jak si spočítat trochu jiný průměr, který vašeho učitele (teda aspoň učitele matiky) přesvědčí, že si tu jedničku zasloužíte.

Začneme pár motivačními příklady.
\EnvId{5lVTVflzYIX-sj4-2bqEx-prumerny-pocet-bodu-dvou-studentu}
\Example{}{
    Dva studenti získali v~testu 70 a~90 bodů. Jaký byl průměrný počet bodů?
}{
    \Answer{$80$ bodů.}

    Spočítáme si
    $$
    \frac{70+90}{2}=80.
    $$
}

\EnvId{TN4dl0yqS7RISW-8ZV4Bu-strana-ctverce-s-obsahem-obdelnika}
\Example{}{
    Obdélník má strany $4\,\rm cm$ a~$9\,\rm cm$. Určete délku strany čtverce, který má stejný obsah jako tento obdélník.
}{
    \Answer{$6\,\rm cm$.}

    Obsah obdélníku je
    $$
    4 \cdot 9 = 36\,\rm cm^2.
    $$
    Abychom dostali stranu čtverce se stejným obsahem, stačí výsledek odmocnit:
    $$
    \sqrt{36}=6.
    $$
    Hledaná strana má tedy délku $6\,\rm cm$.
}

\EnvId{oe3YYvEVymq4SsH63QUnG-prumerna-rychlost-cyklisty}
\Example{}{
    Cyklista ujede:
    \begitems
    \i $60\,\rm km$ rychlostí $30\,\rm km/h$,
    \i $60\,\rm km$ rychlostí $60\,\rm km/h$.
    \enditems
    Jaká je jeho průměrná rychlost?
}{
    \Answer{$40\,\rm km/h$.}

    Průměrnou rychlost dostaneme jako podíl celkové dráhy a~celkového času. Cyklista jel prvních $60\,\rm km$ celkem $2$ hodiny, zbylých $60\,\rm km$ jel $1$ hodinu. Ujel tedy $120\,\rm km$ za $3$ hodiny, a jeho průměrná rychlost je tak
    $$
    40\,\rm km/h.
    $$
    Můžeme to také rozepsat jedním výrazem
    $$
    \frac{120\,\rm km}{\frac{60\,\rm km}{30\,\rm km/h}+\frac{60\,\rm km}{60\,\rm km/h}}
    =
    40\,\rm km/h.
    $$
}

Nyní si jednotlivé průměry představíme trochu víc.

\EnvId{vTih_hh2F2tpJ6pKqYyal-aritmeticky-prumer-dvou-cisel}
\Definition{Aritmetický průměr}{
    Pro dvě reálná čísla $a,b$ definujeme \textbf{aritmetický průměr} jako
    $$
    A=\frac{a+b}{2}.
    $$
}
Aritmetický průměr pro dvě čísla nám dává číslo, co se nachází přešně uprostřed mezi $a$ a~$b$

\EnvId{P1tYaXAyuOdteuzrlQFr3-geometricky-prumer-dvou-cisel}
\Definition{Geometrický průměr}{
    Pro dvě kladná čísla $a,b$ definujeme \textbf{geometrický průměr} jako
    $$
    G=\sqrt{ab}.
    $$
}

Geometrický průměr dvou čísel si můžeme představit jako stranu čtverce, který má stejný obsah jako obdélník se stranami $a$ a~$b$.

\EnvId{ThVI_lBO0FsGdivT12dSO-harmonicky-prumer-dvou-cisel}
\Definition{Harmonický průměr}{
    Pro dvě kladná čísla $a,b$ definujeme \textbf{harmonický průměr} jako
    $$
    H=\frac{2}{\frac1a+\frac1b}=\frac{2ab}{a+b}.
    $$
}

\EnvId{6b6v-mwMDJQLun5LOLiFd-vzorecek-v-prikladu-s-cyklistou}
\Exercise{}{
    Rozmysli si, že jsme použili přesně tento vzoreček v~motivačním příkladu s~cyklistou.
}{
    Máme
    $$
    \frac{120\,\rm km}{60\left(\frac{1}{30\,\rm km/h}+\frac{1}{60\,\rm km/h}\right)}
    =
    \frac{2}{\frac1{30}+\frac1{60}}
    =
    \frac{2 \cdot 30 \cdot 60}{30+60}
    =
    \frac{3600}{90}
    =
    40\,\rm km/h.
    $$
    Použili jsme tedy skutečně harmonický průměr rychlostí $30$ a~$60$.
}

\EnvId{80_Y1Hb5i5TDzLC-gJWHM-porovnani-prumeru-dvojic-cisel}
\Exercise{}{
    Pro následující dvojice čísel spočítej:
    \begitems
    \i aritmetický průměr,
    \i geometrický průměr,
    \i harmonický průměr
    \enditems
    a~je seřaď od nejmenšího po největší:
    $$
    (2,8), \qquad (3,12), \qquad (5,20).
    $$
}{
    \Answer{Pro $(2,8)$: $A=5$, $G=4$, $H=3{,}2$; pro $(3,12)$: $A=7{,}5$, $G=6$, $H=4{,}8$; pro $(5,20)$: $A=12{,}5$, $G=10$, $H=8$. Vždy $H<G<A$.}

    Pro $(2,8)$ dostaneme
    $$
    A=5,\qquad G=4,\qquad H=3{,}2,
    $$
    takže
    $$
    H<G<A.
    $$

    Pro $(3,12)$ dostaneme
    $$
    A=7{,}5,\qquad G=6,\qquad H=4{,}8,
    $$
    a
    $$
    H<G<A.
    $$

    Pro $(5,20)$ dostaneme
    $$
    A=12{,}5,\qquad G=10,\qquad H=8,
    $$
    a zase
    $$
    H<G<A.
    $$
}

Všimneme si, že ve všech případech vyšel aritmetický průměr jako největší, potom geometrický průměr a~harmonický průměr 
jako nejmenší. Bude to platit vždycky? Prozradím, že ano, platí následující nerovnost
$$
H \le G \le A.
$$
Na tobě je tedy najít příklad, ve kterém se všechny průměry budou rovnat.

\EnvId{eisbAdJroxGl2QHbiIPvP-dvojice-s-rovnymi-prumery}
\Exercise{}{
    Najdi dvojici čísel $(a,b)$, pro kterou se jejich aritmetický, geometrický a~harmonický průměr rovnají.
}{
    \Answer{Právě tehdy, když $a=b$.}

    Platí to právě tehdy, když $a=b$.

    Pak totiž
    $$
    A=\frac{a+b}{2}=\frac{a+a}{2}=a,
    $$
    $$
    G=\sqrt{ab}=\sqrt{a^2}=a,
    $$
    $$
    H=\frac{2ab}{a+b}=\frac{2a^2}{2a}=a.
    $$
}

\sec Důkazy nerovností mezi průměry

\EnvId{dyAmE3MykaxHvaQu_8SBc-porovnani-aritmetickeho-a-geometrickeho-prumeru}
\Problem{0}{}{
    Ukaž, že aritmetický průměr dvou čísel je vždy větší nebo roven geometrickému průměru, tedy že pro všechna kladná $a,b$ platí
    $$
    \frac{a+b}{2}\ge \sqrt{ab}.
    $$
}{
    Na minulých hodinách jsme si ukazovali, že $(a-b)^2 \ge 0$, a~z~toho plyne, že $a^2+b^2 \ge 2ab$. Zkus si to připomenout.
}{
    Chceme ukázat, že platí
    $$
    \frac{a+b}{2}\ge \sqrt{ab}
    $$
    protože obě strany jsou nezáporné, můžeme nerovnost umocnit na
    $$
    \left(\frac{a+b}{2}\right)^2 \ge ab.
    $$
    Po úpravě dostaneme
    $$
    \frac{(a+b)^2}{4}\ge ab
    $$
    a vynásobíme čtyřmi a dostaneme
    $$
    (a+b)^2 \ge 4ab.
    $$
    Rozepíšeme levou stranu
    $$
    a^2+2ab+b^2 \ge 4ab
    $$
    a odečteme $4ab$
    $$
    a^2-2ab+b^2 \ge 0.
    $$
    Všimneme si, že levou stranu můžeme přepsat na 
    $$
    (a-b)^2 \ge 0,
    $$
    což platí vždycky, protože druhé mocniny nemůžou být záporné. Tím jsme dokázali danou nerovnost.
}

\EnvId{UV7IZNXLW_o8oeiE9t5jA-porovnani-geometrickeho-a-harmonickeho-prumeru}
\Problem{0}{}{
    Ukaž, že geometrický průměr dvou čísel je vždy větší nebo roven harmonickému průměru, tedy že pro všechna kladná $a,b$ platí
    $$
    \sqrt{ab}\ge \frac{2ab}{a+b}.
    $$
}{
    Zkus postupovat podobně jako v~předchozím důkazu -- zbav se odmocniny a~uprav nerovnost tak, abys dostal tvar $(a-b)^2 \ge 0$.
}{
    Začneme s~nerovností
    $$
    \sqrt{ab}\ge \frac{2ab}{a+b}.
    $$
    Protože $a,b>0$, můžeme ji vynásobit výrazem $a+b$:
    $$
    (a+b)\sqrt{ab}\ge 2ab.
    $$
    Obě strany jsou nezáporné, takže můžeme umocnit a dostaneme
    $$
    (a+b)^2ab \ge 4a^2b^2.
    $$
    Vydělíme výrazem $ab$
    $$
    (a+b)^2 \ge 4ab.
    $$
    a rozepíšeme levou stranu:
    $$
    a^2+2ab+b^2 \ge 4ab.
    $$
    Odečteme $4ab$
    $$
    a^2-2ab+b^2 \ge 0.
    $$
    Teď je evá strana právě
    $$
    (a-b)^2 \ge 0,
    $$
    což už víme, že platí, takže jsme hotovi
}

\sec Definice pro více čísel

Teď se vrátíme ke známkám. Zatím jsme si ukazovali průměry pouze pro dvě čísla, ale můžeme podobně počítat i~pro $n$ čísel, i~když už nevypadají tak pěkně jako předtím a~není nutně vidět, jak tyto průměry vznikají.

\EnvId{aBXQGb9XHMlf2UYpgn2ap-aritmeticky-prumer-n-cisel}
\Definition{Aritmetický průměr ($n$ čísel)}{
    Pro reálná čísla $a_1,a_2,\dots,a_n$ definujeme \textbf{aritmetický průměr} jako
    $$
    A=\frac{a_1+a_2+\cdots+a_n}{n}.
    $$
}

\EnvId{2vR0LCy8VdNIKmBN1e7Vr-geometricky-prumer-n-cisel}
\Definition{Geometrický průměr ($n$ čísel)}{
    Pro kladná čísla $a_1,a_2,\dots,a_n$ definujeme \textbf{geometrický průměr} jako
    $$
    G=\root n \of {a_1 a_2 \cdots a_n}.
    $$
}

\EnvId{RkjiONhHfMoT3_KECSx68-harmonicky-prumer-n-cisel}
\Definition{Harmonický průměr ($n$ čísel)}{
    Pro kladná čísla $a_1,a_2,\dots,a_n$ definujeme \textbf{harmonický průměr} jako
    $$
    H=\frac{n}{\frac1{a_1}+\frac1{a_2}+\cdots+\frac1{a_n}}.
    $$
}

\sec Známky

\EnvId{vUtBaQxWMNYAPcJPGjAhI-vypocet-prumeru-znamek}
\Example{}{
    Pomocí těchto vzorečků spočítejte aritmetický, geometrický a~harmonický průměr známek
    $$
    1,\ 1,\ 3,\ 2,\ 1
    $$
    (na geometrický průměr použijte kalkulačku).
}{
    \Answer{$A=1{,}6$, $G\approx 1{,}43$, $H\approx 1{,}30$.}

    Aritmetický průměr:
    $$
    A=\frac{1+1+3+2+1}{5}=1{,}6.
    $$

    Geometrický průměr:
    $$
    G=\root 5 \of {1 \cdot 1 \cdot 3 \cdot 2 \cdot 1}=\root 5 \of 6 \approx 1{,}43.
    $$

    Harmonický průměr:
    $$
    H=\frac{5}{1+1+\frac13+\frac12+1}\approx 1{,}30.
    $$
}

Všimni si, že zase platí, že aritmetický průměr je největší, potom geometrický a~harmonický je nejmenší. Dokázali jsme to pouze pro dvojici čísel, ale neostré nerovnosti mezi průměry platí pro všechna $n$, jen jejich důkazy jsou mnohem složitější.

Ale ukázali jsme si, že když ti bude vycházet horší známka na vysvědčení, tak si můžeš spočítat harmonický průměr a~pokusit se přesvědčit učitele, že ten se zaokrouhlí nahoru.

\DisplayExerciseSolutions

\DisplayProblemSolutions

\bye